\documentclass[%
 reprint,
%superscriptaddress,
%groupedaddress,
%unsortedaddress,
%runinaddress,
%frontmatterverbose, 
%preprint,
%preprintnumbers,
%nofootinbib,
%nobibnotes,
%bibnotes,
 amsmath,amssymb,
 aps,
%pra,
%prb,
%rmp,
%prstab,
%prstper,
%floatfix,
]{revtex4-2}

\usepackage{graphicx}
\usepackage{dcolumn}
\usepackage{bm}
\usepackage{kotex}
\usepackage{hyperref}

\usepackage{braket}

%\usepackage[mathlines]{lineno}% Enable numbering of text and display math
%\linenumbers\relax % Commence numbering lines

%\usepackage[showframe,%Uncomment any one of the following lines to test 
%%scale=0.7, marginratio={1:1, 2:3}, ignoreall,% default settings
%%text={7in,10in},centering,
%%margin=1.5in,
%%total={6.5in,8.75in}, top=1.2in, left=0.9in, includefoot,
%%height=10in,a5paper,hmargin={3cm,0.8in},
%]{geometry}

\begin{document}

\preprint{APS/123-QED}

\title{Hall Effect}
% \thanks{A footnote to the article title}

\author{박찬우}
\email{pco0511@kaist.ac.kr}
\author{최유승}
\email{yooseung1121@kaist.ac.kr}
\author{임원준}
\email{wonjunlim@kaist.ac.kr}
\affiliation{Physics Department, KAIST.}

\date{\today}

\begin{abstract}
An article usually includes an abstract, a concise summary of the work
covered at length in the main body of the article. 
\end{abstract}

\maketitle

%\tableofcontents

\section{\label{sec:intro}Introduction}

intro placeholder

\section{\label{sec:theory}Theoretical background}

theory placeholder

\section{\label{sec:methods}Experimental Methods}

methods placeholder

\section{\label{sec:results}Results}

results placeholder

\section{\label{sec:discussion}Discussion}

discussion placeholder

\section{\label{sec:conclusion}Conclusion}

conclusion placeholder





% \begin{table}[b]%The best place to locate the table environment is directly after its first reference in text
% \caption{\label{tab:table1}%
% A table that fits into a single column of a two-column layout. 
% Note that REV\TeX~4 adjusts the intercolumn spacing so that the table fills the
% entire width of the column. Table captions are numbered
% automatically. 
% This table illustrates left-, center-, decimal- and right-aligned columns,
% along with the use of the \texttt{ruledtabular} environment which sets the 
% Scotch (double) rules above and below the alignment, per APS style.
% }
% \begin{ruledtabular}
% \begin{tabular}{lcdr}
% \textrm{Left\footnote{Note a.}}&
% \textrm{Centered\footnote{Note b.}}&
% \multicolumn{1}{c}{\textrm{Decimal}}&
% \textrm{Right}\\
% \colrule
% 1 & 2 & 3.001 & 4\\
% 10 & 20 & 30 & 40\\
% 100 & 200 & 300.0 & 400\\
% \end{tabular}
% \end{ruledtabular}
% \end{table}
% and Fig.~\ref{fig:epsart}.%
% \begin{figure}[b]
% \includegraphics{fig_1}% Here is how to import EPS art
% \caption{\label{fig:epsart} A figure caption. The figure captions are
% automatically numbered.}
% \end{figure}



% \begin{figure*}
% \includegraphics{fig_2}% Here is how to import EPS art
% \caption{\label{fig:wide}Use the figure* environment to get a wide
% figure that spans the page in \texttt{twocolumn} formatting.}
% \end{figure*}


% \begin{table*}
% \caption{\label{tab:table3}This is a wide table that spans the full page
% width in a two-column layout. It is formatted using the
% \texttt{table*} environment. It also demonstates the use of
% \textbackslash\texttt{multicolumn} in rows with entries that span
% more than one column.}
% \begin{ruledtabular}
% \begin{tabular}{ccccc}
%  &\multicolumn{2}{c}{$D_{4h}^1$}&\multicolumn{2}{c}{$D_{4h}^5$}\\
%  Ion&1st alternative&2nd alternative&lst alternative
% &2nd alternative\\ \hline
%  K&$(2e)+(2f)$&$(4i)$ &$(2c)+(2d)$&$(4f)$ \\
%  Mn&$(2g)$\footnote{The $z$ parameter of these positions is $z\sim\frac{1}{4}$.}
%  &$(a)+(b)+(c)+(d)$&$(4e)$&$(2a)+(2b)$\\
%  Cl&$(a)+(b)+(c)+(d)$&$(2g)$\footnotemark[1]
%  &$(4e)^{\text{a}}$\\
%  He&$(8r)^{\text{a}}$&$(4j)^{\text{a}}$&$(4g)^{\text{a}}$\\
%  Ag& &$(4k)^{\text{a}}$& &$(4h)^{\text{a}}$\\
% \end{tabular}
% \end{ruledtabular}
% \end{table*}


% \begin{table}[b]
% \caption{\label{tab:table4}%
% Numbers in columns Three--Five are aligned with the ``d'' column specifier 
% (requires the \texttt{dcolumn} package). 
% Non-numeric entries (those entries without a ``.'') in a ``d'' column are aligned on the decimal point. 
% Use the ``D'' specifier for more complex layouts. }
% \begin{ruledtabular}
% \begin{tabular}{ccddd}
% One&Two&
% \multicolumn{1}{c}{\textrm{Three}}&
% \multicolumn{1}{c}{\textrm{Four}}&
% \multicolumn{1}{c}{\textrm{Five}}\\
% %\mbox{Three}&\mbox{Four}&\mbox{Five}\\
% \hline
% one&two&\mbox{three}&\mbox{four}&\mbox{five}\\
% He&2& 2.77234 & 45672. & 0.69 \\
% C\footnote{Some tables require footnotes.}
%   &C\footnote{Some tables need more than one footnote.}
%   & 12537.64 & 37.66345 & 86.37 \\
% \end{tabular}
% \end{ruledtabular}
% \end{table}


% Tables~\ref{tab:table1}, \ref{tab:table3}, \ref{tab:table4}, and \ref{tab:table2}%
% \begin{table}[b]
% \caption{\label{tab:table2}
% A table with numerous columns that still fits into a single column. 
% Here, several entries share the same footnote. 
% Inspect the \LaTeX\ input for this table to see exactly how it is done.}
% \begin{ruledtabular}
% \begin{tabular}{cccccccc}
%  &$r_c$ (\AA)&$r_0$ (\AA)&$\kappa r_0$&
%  &$r_c$ (\AA) &$r_0$ (\AA)&$\kappa r_0$\\
% \hline
% Cu& 0.800 & 14.10 & 2.550 &Sn\footnotemark[1]
% & 0.680 & 1.870 & 3.700 \\
% Ag& 0.990 & 15.90 & 2.710 &Pb\footnotemark[2]
% & 0.450 & 1.930 & 3.760 \\
% Au& 1.150 & 15.90 & 2.710 &Ca\footnotemark[3]
% & 0.750 & 2.170 & 3.560 \\
% Mg& 0.490 & 17.60 & 3.200 &Sr\footnotemark[4]
% & 0.900 & 2.370 & 3.720 \\
% Zn& 0.300 & 15.20 & 2.970 &Li\footnotemark[2]
% & 0.380 & 1.730 & 2.830 \\
% Cd& 0.530 & 17.10 & 3.160 &Na\footnotemark[5]
% & 0.760 & 2.110 & 3.120 \\
% Hg& 0.550 & 17.80 & 3.220 &K\footnotemark[5]
% &  1.120 & 2.620 & 3.480 \\
% Al& 0.230 & 15.80 & 3.240 &Rb\footnotemark[3]
% & 1.330 & 2.800 & 3.590 \\
% Ga& 0.310 & 16.70 & 3.330 &Cs\footnotemark[4]
% & 1.420 & 3.030 & 3.740 \\
% In& 0.460 & 18.40 & 3.500 &Ba\footnotemark[5]
% & 0.960 & 2.460 & 3.780 \\
% Tl& 0.480 & 18.90 & 3.550 & & & & \\
% \end{tabular}
% \end{ruledtabular}
% \footnotetext[1]{Here's the first, from Ref.~\onlinecite{feyn54}.}
% \footnotetext[2]{Here's the second.}
% \footnotetext[3]{Here's the third.}
% \footnotetext[4]{Here's the fourth.}
% \footnotetext[5]{And etc.}
% \end{table}



\appendix

\section{Appendixes}

appendix placeholder

% \nocite{*}

\bibliography{main}

\end{document}